
\documentclass{article}
\usepackage{colortbl}
\usepackage{makecell}
\usepackage{multirow}
\usepackage{supertabular}

\begin{document}

\newcounter{utterance}

\centering \large Interaction Transcript for game `guesswhat', experiment `Abs\_Level\_2', episode 7 with reflexion{-}llama{-}3.1{-}8b{-}instruct.
\vspace{24pt}

{ \footnotesize  \setcounter{utterance}{1}
\setlength{\tabcolsep}{0pt}
\begin{supertabular}{c@{$\;$}|p{.15\linewidth}@{}p{.15\linewidth}p{.15\linewidth}p{.15\linewidth}p{.15\linewidth}p{.15\linewidth}}
    \# & $\;$A & \multicolumn{4}{c}{Game Master} & $\;\:$B\\
    \hline

    \theutterance \stepcounter{utterance}  
    & & \multicolumn{4}{p{0.6\linewidth}}{
        \cellcolor[rgb]{0.9,0.9,0.9}{
            \makecell[{{p{\linewidth}}}]{
                \texttt{\tiny{[P1$\langle$GM]}}
                \texttt{You have attempted this same task before and did not succeed. These are the plans you wrote after those attempts.} \\
\texttt{Trial \#1: New plan: I got stuck in a loop of asking yes/no questions that were too specific and didn't help me narrow down the options effectively. I also didn't use the information I gathered from previous questions to inform my next questions. Next time, I will ask more general questions that can help me eliminate a larger number of options at once. For example, I could ask if the target word is a noun, or if it's something that can be found indoors. I will also try to use the information I gather from previous questions to inform my next questions, and avoid asking questions that are too similar to previous ones.} \\
\\ 
\texttt{Trial \#2: New plan: In my previous attempt, I got stuck in a loop of asking yes/no questions that were too specific and didn't help me narrow down the options effectively. I also didn't use the information I gathered from previous questions to inform my next questions. However, I did manage to gather some useful information, such as the fact that the target word is man{-}made, can be found indoors, and is a geometric shape. } \\
\\ 
\texttt{To improve my strategy, I will ask more general questions that can help me eliminate a larger number of options at once. For example, I could ask if the target word is a type of polygon, or if it's a shape that can be found in a specific room or area of the house. I will also try to use the information I gather from previous questions to inform my next questions, and avoid asking questions that are too similar to previous ones. } \\
\\ 
\texttt{Additionally, I will pay closer attention to the answers I receive and try to make connections between them. For example, if I ask if the target word is a type of polygon and the answer is yes, I will then ask more specific questions about the number of sides or angles of the polygon. } \\
\\ 
\texttt{My next question will be: QUESTION: Is the target word a type of polygon?} \\
\\ 
\texttt{Trial \#3: New plan: In my previous attempts, I got stuck in a loop of asking yes/no questions that were too specific and didn't help me narrow down the options effectively. I also didn't use the information I gathered from previous questions to inform my next questions. However, I did manage to gather some useful information, such as the fact that the target word is man{-}made, can be found indoors, and is a geometric shape.} \\
\\ 
\texttt{To improve my strategy, I will ask more general questions that can help me eliminate a larger number of options at once. For example, I could ask if the target word is related to architecture, or if it's a shape that can be found in a specific room or area of the house. I will also try to use the information I gather from previous questions to inform my next questions, and avoid asking questions that are too similar to previous ones.} \\
\\ 
\texttt{Additionally, I will pay closer attention to the answers I receive and try to make connections between them. For example, if I ask if the target word is related to architecture and the answer is yes, I will then ask more specific questions about the type of architecture or the materials used to build it.} \\
\\ 
\texttt{I will also try to ask more open{-}ended questions that can help me gather more information about the target word. For example, I could ask if the target word is something that can be held in one's hand, or if it's a shape that can be found in nature.} \\
\\ 
\texttt{My next question will be: QUESTION: Is the target word related to architecture?} \\
\\ 
\\ 
\texttt{Let's play a guessing game. Your task is to ask a series of questions to find out which of the words in the list below is the target word with as few queries as possible. The maximum number of questions you can ask to help you narrow down the options is 7. Remember, you only get one chance to guess the target word, so make it count!} \\
\\ 
\texttt{The rules of the game are: } \\
\\ 
\texttt{1. Only ask one question per turn. DO NOT ASK MORE THAN ONE QUESTION.} \\
\\ 
\texttt{2. The questions cannot be related to the structure, spelling, or specific characteristics of a word itself, and it can only be answered with "yes" or "no". } \\
\\ 
\texttt{3. Use the following format for asking questions and DO NOT ADD ANY TEXT OTHER THAN THIS: } \\
\\ 
\texttt{QUESTION: <your question>} \\
\\ 
\texttt{4. You have only one opportunity to guess the target word. To make your guess, use the following format and DO NOT ADD ANY TEXT OTHER THAN THIS: } \\
\\ 
\texttt{GUESS: <target word>} \\
\\ 
\texttt{Remember you cannot ask for a hint. Make your question brief. Let’s start!} \\
\\ 
\texttt{List of candidate words: {[}'grandson', 'granddaughter', 'mother{-}in{-}law', 'sister{-}in{-}law', 'trapezium', 'parallelogram', 'decahedron', 'dodecahedron'{]}} \\
            }
        }
    }
    & \\ \\

    \theutterance \stepcounter{utterance}  
    & \multicolumn{4}{p{0.6\linewidth}}{
        \cellcolor[rgb]{0.9,0.9,0.9}{
            \makecell[{{p{\linewidth}}}]{
                \texttt{\tiny{[P1$\rangle$GM]}}
                \texttt{QUESTION: Is the target word man{-}made?} \\
            }
        }
    }
    & & \\ \\

    \theutterance \stepcounter{utterance}  
    & & \multicolumn{4}{p{0.6\linewidth}}{
        \cellcolor[rgb]{0.9,0.9,0.9}{
            \makecell[{{p{\linewidth}}}]{
                \texttt{\tiny{[GM$\rangle$P2]}}
                \texttt{You have attempted this same task before and did not succeed. These are the plans you wrote after those attempts.} \\
\texttt{Trial \#1: New plan: I failed to guess the target word because I didn't ask enough questions to narrow down the possibilities. I should have asked more questions about the geometric shape and its properties, such as the number of sides, the type of polygon, and any specific features it may have. I should also have paid closer attention to the previous questions and answers to avoid asking redundant or irrelevant questions. Next time, I will ask more targeted and thoughtful questions to gather more information and increase my chances of guessing the target word correctly.} \\
\\ 
\texttt{Trial \#2: New plan: I failed to guess the target word because I didn't ask enough questions to narrow down the possibilities, and I also didn't pay close attention to the previous questions and answers. However, in this attempt, I was able to gather some useful information, such as the fact that the target word is a geometric shape with more than four sides. Next time, I will build on this information and ask more targeted and thoughtful questions to gather more details about the shape, such as its type (e.g. polygon, polyhedron), its properties (e.g. regular, irregular), and any specific features it may have. I will also pay closer attention to the previous questions and answers to avoid asking redundant or irrelevant questions. Additionally, I will try to ask more questions that can be answered with a "yes" or "no" to get more information without giving away too much.} \\
\\ 
\texttt{Trial \#3: New plan: I failed to guess the target word because I didn't ask enough questions to narrow down the possibilities, and I also didn't pay close attention to the previous questions and answers. However, in this attempt, I was able to gather some useful information, such as the fact that the target word is a geometric shape with more than four sides. Next time, I will build on this information and ask more targeted and thoughtful questions to gather more details about the shape, such as its type (e.g. polygon, polyhedron), its properties (e.g. regular, irregular), and any specific features it may have. I will also pay closer attention to the previous questions and answers to avoid asking redundant or irrelevant questions. Additionally, I will try to ask more questions that can be answered with a "yes" or "no" to get more information without giving away too much. Furthermore, I will also try to ask questions that can help me to eliminate possibilities, such as "Is the target word a type of polyhedron with a specific number of faces?" or "Does the target word have a specific symmetry?" This will help me to narrow down the possibilities and increase my chances of guessing the target word correctly.} \\
\\ 
\\ 
\texttt{Let’s play a guessing game! Your task is to answer the other player's questions. Based on your knowledge of the word: decahedron, respond to the following questions or guesses. Limit your response to only 'yes' or 'no' with no explanation or other words. Never reveal the answer in your response.} \\
\\ 
\texttt{You must reply using the format below and DO NOT ADD ANY TEXT OTHER THAN THIS:} \\
\\ 
\texttt{ANSWER: <some text>} \\
\\ 
\texttt{Target Word: decahedron} \\
\\ 
\texttt{QUESTION: Is the target word man{-}made?} \\
            }
        }
    }
    & \\ \\

    \theutterance \stepcounter{utterance}  
    & & & \multicolumn{4}{p{0.6\linewidth}}{
        \cellcolor[rgb]{0.9,0.9,0.9}{
            \makecell[{{p{\linewidth}}}]{
                \texttt{\tiny{[GM$\langle$P2]}}
                \texttt{ANSWER: yes} \\
            }
        }
    }
     \\ \\

    \theutterance \stepcounter{utterance}  
    & & \multicolumn{4}{p{0.6\linewidth}}{
        \cellcolor[rgb]{0.9,0.9,0.9}{
            \makecell[{{p{\linewidth}}}]{
                \texttt{\tiny{[P1$\langle$GM]}}
                \texttt{ANSWER: yes} \\
            }
        }
    }
    & \\ \\

    \theutterance \stepcounter{utterance}  
    & \multicolumn{4}{p{0.6\linewidth}}{
        \cellcolor[rgb]{0.9,0.9,0.9}{
            \makecell[{{p{\linewidth}}}]{
                \texttt{\tiny{[P1$\rangle$GM]}}
                \texttt{QUESTION: Is the target word a geometric shape?} \\
            }
        }
    }
    & & \\ \\

    \theutterance \stepcounter{utterance}  
    & & \multicolumn{4}{p{0.6\linewidth}}{
        \cellcolor[rgb]{0.9,0.9,0.9}{
            \makecell[{{p{\linewidth}}}]{
                \texttt{\tiny{[GM$\rangle$P2]}}
                \texttt{QUESTION: Is the target word a geometric shape?} \\
            }
        }
    }
    & \\ \\

    \theutterance \stepcounter{utterance}  
    & & & \multicolumn{4}{p{0.6\linewidth}}{
        \cellcolor[rgb]{0.9,0.9,0.9}{
            \makecell[{{p{\linewidth}}}]{
                \texttt{\tiny{[GM$\langle$P2]}}
                \texttt{ANSWER: yes} \\
            }
        }
    }
     \\ \\

    \theutterance \stepcounter{utterance}  
    & & \multicolumn{4}{p{0.6\linewidth}}{
        \cellcolor[rgb]{0.9,0.9,0.9}{
            \makecell[{{p{\linewidth}}}]{
                \texttt{\tiny{[P1$\langle$GM]}}
                \texttt{ANSWER: yes} \\
            }
        }
    }
    & \\ \\

    \theutterance \stepcounter{utterance}  
    & \multicolumn{4}{p{0.6\linewidth}}{
        \cellcolor[rgb]{0.9,0.9,0.9}{
            \makecell[{{p{\linewidth}}}]{
                \texttt{\tiny{[P1$\rangle$GM]}}
                \texttt{QUESTION: Is the target word a polygon?} \\
            }
        }
    }
    & & \\ \\

    \theutterance \stepcounter{utterance}  
    & & \multicolumn{4}{p{0.6\linewidth}}{
        \cellcolor[rgb]{0.9,0.9,0.9}{
            \makecell[{{p{\linewidth}}}]{
                \texttt{\tiny{[GM$\rangle$P2]}}
                \texttt{QUESTION: Is the target word a polygon?} \\
            }
        }
    }
    & \\ \\

    \theutterance \stepcounter{utterance}  
    & & & \multicolumn{4}{p{0.6\linewidth}}{
        \cellcolor[rgb]{0.9,0.9,0.9}{
            \makecell[{{p{\linewidth}}}]{
                \texttt{\tiny{[GM$\langle$P2]}}
                \texttt{ANSWER: no} \\
            }
        }
    }
     \\ \\

    \theutterance \stepcounter{utterance}  
    & & \multicolumn{4}{p{0.6\linewidth}}{
        \cellcolor[rgb]{0.9,0.9,0.9}{
            \makecell[{{p{\linewidth}}}]{
                \texttt{\tiny{[P1$\langle$GM]}}
                \texttt{ANSWER: no} \\
            }
        }
    }
    & \\ \\

    \theutterance \stepcounter{utterance}  
    & \multicolumn{4}{p{0.6\linewidth}}{
        \cellcolor[rgb]{0.9,0.9,0.9}{
            \makecell[{{p{\linewidth}}}]{
                \texttt{\tiny{[P1$\rangle$GM]}}
                \texttt{QUESTION: Is the target word something that can be found indoors?} \\
            }
        }
    }
    & & \\ \\

    \theutterance \stepcounter{utterance}  
    & & \multicolumn{4}{p{0.6\linewidth}}{
        \cellcolor[rgb]{0.9,0.9,0.9}{
            \makecell[{{p{\linewidth}}}]{
                \texttt{\tiny{[GM$\rangle$P2]}}
                \texttt{QUESTION: Is the target word something that can be found indoors?} \\
            }
        }
    }
    & \\ \\

    \theutterance \stepcounter{utterance}  
    & & & \multicolumn{4}{p{0.6\linewidth}}{
        \cellcolor[rgb]{0.9,0.9,0.9}{
            \makecell[{{p{\linewidth}}}]{
                \texttt{\tiny{[GM$\langle$P2]}}
                \texttt{ANSWER: yes} \\
            }
        }
    }
     \\ \\

    \theutterance \stepcounter{utterance}  
    & & \multicolumn{4}{p{0.6\linewidth}}{
        \cellcolor[rgb]{0.9,0.9,0.9}{
            \makecell[{{p{\linewidth}}}]{
                \texttt{\tiny{[P1$\langle$GM]}}
                \texttt{ANSWER: yes} \\
            }
        }
    }
    & \\ \\

    \theutterance \stepcounter{utterance}  
    & \multicolumn{4}{p{0.6\linewidth}}{
        \cellcolor[rgb]{0.9,0.9,0.9}{
            \makecell[{{p{\linewidth}}}]{
                \texttt{\tiny{[P1$\rangle$GM]}}
                \texttt{QUESTION: Is the target word a type of polyhedron?} \\
            }
        }
    }
    & & \\ \\

    \theutterance \stepcounter{utterance}  
    & & \multicolumn{4}{p{0.6\linewidth}}{
        \cellcolor[rgb]{0.9,0.9,0.9}{
            \makecell[{{p{\linewidth}}}]{
                \texttt{\tiny{[GM$\rangle$P2]}}
                \texttt{QUESTION: Is the target word a type of polyhedron?} \\
            }
        }
    }
    & \\ \\

    \theutterance \stepcounter{utterance}  
    & & & \multicolumn{4}{p{0.6\linewidth}}{
        \cellcolor[rgb]{0.9,0.9,0.9}{
            \makecell[{{p{\linewidth}}}]{
                \texttt{\tiny{[GM$\langle$P2]}}
                \texttt{ANSWER: yes} \\
            }
        }
    }
     \\ \\

    \theutterance \stepcounter{utterance}  
    & & \multicolumn{4}{p{0.6\linewidth}}{
        \cellcolor[rgb]{0.9,0.9,0.9}{
            \makecell[{{p{\linewidth}}}]{
                \texttt{\tiny{[P1$\langle$GM]}}
                \texttt{ANSWER: yes} \\
            }
        }
    }
    & \\ \\

    \theutterance \stepcounter{utterance}  
    & \multicolumn{4}{p{0.6\linewidth}}{
        \cellcolor[rgb]{0.9,0.9,0.9}{
            \makecell[{{p{\linewidth}}}]{
                \texttt{\tiny{[P1$\rangle$GM]}}
                \texttt{QUESTION: Does the target word have more than 4 faces?} \\
            }
        }
    }
    & & \\ \\

    \theutterance \stepcounter{utterance}  
    & & \multicolumn{4}{p{0.6\linewidth}}{
        \cellcolor[rgb]{0.9,0.9,0.9}{
            \makecell[{{p{\linewidth}}}]{
                \texttt{\tiny{[GM$\rangle$P2]}}
                \texttt{QUESTION: Does the target word have more than 4 faces?} \\
            }
        }
    }
    & \\ \\

    \theutterance \stepcounter{utterance}  
    & & & \multicolumn{4}{p{0.6\linewidth}}{
        \cellcolor[rgb]{0.9,0.9,0.9}{
            \makecell[{{p{\linewidth}}}]{
                \texttt{\tiny{[GM$\langle$P2]}}
                \texttt{ANSWER: yes} \\
            }
        }
    }
     \\ \\

    \theutterance \stepcounter{utterance}  
    & & \multicolumn{4}{p{0.6\linewidth}}{
        \cellcolor[rgb]{0.9,0.9,0.9}{
            \makecell[{{p{\linewidth}}}]{
                \texttt{\tiny{[P1$\langle$GM]}}
                \texttt{ANSWER: yes} \\
            }
        }
    }
    & \\ \\

    \theutterance \stepcounter{utterance}  
    & \multicolumn{4}{p{0.6\linewidth}}{
        \cellcolor[rgb]{0.9,0.9,0.9}{
            \makecell[{{p{\linewidth}}}]{
                \texttt{\tiny{[P1$\rangle$GM]}}
                \texttt{QUESTION: Does the target word have 12 faces?} \\
            }
        }
    }
    & & \\ \\

    \theutterance \stepcounter{utterance}  
    & & \multicolumn{4}{p{0.6\linewidth}}{
        \cellcolor[rgb]{0.9,0.9,0.9}{
            \makecell[{{p{\linewidth}}}]{
                \texttt{\tiny{[GM$\rangle$P2]}}
                \texttt{QUESTION: Does the target word have 12 faces?} \\
            }
        }
    }
    & \\ \\

    \theutterance \stepcounter{utterance}  
    & & & \multicolumn{4}{p{0.6\linewidth}}{
        \cellcolor[rgb]{0.9,0.9,0.9}{
            \makecell[{{p{\linewidth}}}]{
                \texttt{\tiny{[GM$\langle$P2]}}
                \texttt{ANSWER: yes} \\
            }
        }
    }
     \\ \\

    \theutterance \stepcounter{utterance}  
    & & \multicolumn{4}{p{0.6\linewidth}}{
        \cellcolor[rgb]{0.9,0.9,0.9}{
            \makecell[{{p{\linewidth}}}]{
                \texttt{\tiny{[P1$\langle$GM]}}
                \texttt{ANSWER: yes} \\
            }
        }
    }
    & \\ \\

    \theutterance \stepcounter{utterance}  
    & \multicolumn{4}{p{0.6\linewidth}}{
        \cellcolor[rgb]{0.9,0.9,0.9}{
            \makecell[{{p{\linewidth}}}]{
                \texttt{\tiny{[P1$\rangle$GM]}}
                \texttt{GUESS: dodecahedron} \\
            }
        }
    }
    & & \\ \\

    \theutterance \stepcounter{utterance}  
    & & & \multicolumn{2}{p{0.3\linewidth}}{
        \cellcolor[rgb]{0.9,0.9,0.9}{
            \makecell[{{p{\linewidth}}}]{
                \texttt{\tiny{[GM$|$GM]}}
                \texttt{dodecahedron} \\
            }
        }
    }
    & & \\ \\

    \theutterance \stepcounter{utterance}  
    & & & \multicolumn{2}{p{0.3\linewidth}}{
        \cellcolor[rgb]{0.9,0.9,0.9}{
            \makecell[{{p{\linewidth}}}]{
                \texttt{\tiny{[GM$|$GM]}}
                \texttt{end game} \\
            }
        }
    }
    & & \\ \\

    \theutterance \stepcounter{utterance}  
    & \multicolumn{4}{p{0.6\linewidth}}{
        \cellcolor[rgb]{0.95,0.95,0.95}{
            \makecell[{{p{\linewidth}}}]{
                \texttt{\tiny{[P2$\rightarrow$P1]}}
                \texttt{Post{-}episode Reflexion (not sent to GM):} \\
\texttt{New plan: In my previous attempts, I got stuck in a loop of asking yes/no questions that were too specific and didn't help me narrow down the options effectively. I also didn't use the information I gathered from previous questions to inform my next questions. However, I did manage to gather some useful information, such as the fact that the target word is man{-}made, can be found indoors, and is a geometric shape.} \\
\\ 
\texttt{To improve my strategy, I will ask more general questions that can help me eliminate a larger number of options at once. For example, I could ask if the target word is related to a specific type of object or category, such as furniture, decorations, or tools. I will also try to use the information I gather from previous questions to inform my next questions, and avoid asking questions that are too similar to previous ones.} \\
\\ 
\texttt{Additionally, I will pay closer attention to the answers I receive and try to make connections between them. For example, if I ask if the target word is related to a specific type of object and the answer is yes, I will then ask more specific questions about the characteristics or features of that object.} \\
\\ 
\texttt{I will also try to ask more open{-}ended questions that can help me gather more information about the target word. For example, I could ask if the target word is something that serves a functional purpose, or if it's a shape that can be found in a specific room or area of the house.} \\
\\ 
\texttt{My next question will be: QUESTION: Is the target word something that serves a functional purpose?} \\
            }
        }
    }
    & & \\ \\

    \theutterance \stepcounter{utterance}  
    & & & \multicolumn{4}{p{0.6\linewidth}}{
        \cellcolor[rgb]{0.95,0.95,0.95}{
            \makecell[{{p{\linewidth}}}]{
                \texttt{\tiny{[P1$\rightarrow$P2]}}
                \texttt{Post{-}episode Reflexion (not sent to GM):} \\
\texttt{New plan: I failed to guess the target word because I didn't ask enough questions to narrow down the possibilities, and I also didn't pay close attention to the previous questions and answers. However, in this attempt, I was able to gather some useful information, such as the fact that the target word is a geometric shape with more than four sides, is a type of polyhedron, and has 12 faces. Next time, I will build on this information and ask more targeted and thoughtful questions to gather more details about the shape, such as its properties (e.g. regular, irregular), and any specific features it may have. I will also pay closer attention to the previous questions and answers to avoid asking redundant or irrelevant questions. Additionally, I will try to ask more questions that can be answered with a "yes" or "no" to get more information without giving away too much. Furthermore, I will also try to ask questions that can help me to eliminate possibilities, such as "Is the target word a regular polyhedron?" or "Does the target word have a specific type of symmetry?" This will help me to narrow down the possibilities and increase my chances of guessing the target word correctly.} \\
            }
        }
    }
     \\ \\

\end{supertabular}
}

\end{document}
